%! Rates of Change in Linear and Quadratic Functions — Unit 1, Lesson 1.3 — Group Activity
\documentclass[10pt]{article}
\usepackage{precalculus-article}
\usepackage{precalculus-boxes}
\newcommand{\TallMath}[1]{$\displaystyle #1\rule[-1.4em]{0pt}{3.2em}$}

\begin{document}

\pageheader{Unit 1, Lesson 1.3}{Group Activity}
\namepartnerperiod

\vspace{0.12in}

\begin{scenariobox}[A Ball in Free Fall]{sky}
A ball is dropped from a height of 144 feet. Its height above the ground (in feet) is given by
$h(t) = 144 - 16t^2$, where $t$ is time in seconds. The table below shows heights at half-second intervals.

\smallskip
\renewcommand{\arraystretch}{1.4}
\begin{tabular}{|c|c|c|c|c|c|c|c|}
\hline
$t$ (sec)   & 0   & 0.5 & 1.0 & 1.5 & 2.0  & 2.5  & 3.0 \\
\hline
$h(t)$ (ft) & 144 & 140 & 128 & 108 & 80   & 44   & 0   \\
\hline
\end{tabular}
\renewcommand{\arraystretch}{1}
\end{scenariobox}

\vspace{0.1in}

\begin{tcolorbox}[colback=white, colframe=black!40,
    title=\textbf{Tier R --- Remediate},
    fonttitle=\bfseries, arc=2mm, left=3mm, right=3mm, top=2mm, bottom=2mm]

\textbf{Directions:} Compute the AROC of height over each 0.5-second interval.

\smallskip
\renewcommand{\arraystretch}{1.5}
\begin{tabular}{|l|c|}
\hline
\textbf{Interval} & \textbf{AROC (ft/sec)} \\
\hline
$[0,\; 0.5]$    & \blank{3cm} \\
$[0.5,\; 1.0]$  & \blank{3cm} \\
$[1.0,\; 1.5]$  & \blank{3cm} \\
$[1.5,\; 2.0]$  & \blank{3cm} \\
$[2.0,\; 2.5]$  & \blank{3cm} \\
$[2.5,\; 3.0]$  & \blank{3cm} \\
\hline
\end{tabular}
\renewcommand{\arraystretch}{1}

\vspace{0.1in}
Are the AROCs increasing or decreasing? \blank{4cm}

\vspace{0.08in}
What does this tell you about the ball's speed as it falls? \writeline \writeline
\end{tcolorbox}

\vspace{0.1in}

\begin{tcolorbox}[colback=white, colframe=black!40,
    title=\textbf{Tier A --- Approaching Proficiency},
    fonttitle=\bfseries, arc=2mm, left=3mm, right=3mm, top=2mm, bottom=2mm]

\textbf{Directions:} Using the AROCs from Tier R, compute the second differences.

\smallskip
\renewcommand{\arraystretch}{1.5}
\begin{tabular}{|l|c|c|}
\hline
\textbf{Interval} & \textbf{AROC (first diff.)} & \textbf{Change in AROC (second diff.)} \\
\hline
$[0,\; 0.5]$    & \blank{2.5cm} & \\
$[0.5,\; 1.0]$  & \blank{2.5cm} & \blank{2.5cm} \\
$[1.0,\; 1.5]$  & \blank{2.5cm} & \blank{2.5cm} \\
$[1.5,\; 2.0]$  & \blank{2.5cm} & \blank{2.5cm} \\
$[2.0,\; 2.5]$  & \blank{2.5cm} & \blank{2.5cm} \\
$[2.5,\; 3.0]$  & \blank{2.5cm} & \blank{2.5cm} \\
\hline
\end{tabular}
\renewcommand{\arraystretch}{1}

\vspace{0.1in}
Are the second differences constant? \blank{2cm}

\vspace{0.08in}
What function type does this constant second difference confirm for $h(t)$? \writeline
\end{tcolorbox}

\vspace{0.1in}

\begin{tcolorbox}[colback=white, colframe=black!40,
    title=\textbf{Tier E --- Extension},
    fonttitle=\bfseries, arc=2mm, left=3mm, right=3mm, top=2mm, bottom=2mm]

\textbf{Directions:} Use the constant second difference you found to \textit{predict} the AROC on $[3.0, 3.5]$ without computing from the formula. Then verify.

\vspace{0.06in}
\textbf{Prediction:} The AROC on $[2.5, 3.0]$ is \blank{2cm}. The constant second difference is \blank{2cm}.

Therefore, I predict the AROC on $[3.0, 3.5]$ is \blank{3cm}.

\vspace{0.1in}
\textbf{Verification:} Compute $h(3.5)$ using the formula and find the actual AROC on $[3.0, 3.5]$.

\vspace{0.5in}
Actual AROC: \blank{3cm} \quad Does it match your prediction? \blank{2cm}

\vspace{0.1in}
\textbf{Physics connection:} In words, explain why the second differences for a freely falling
object are constant. (\textit{Hint: What quantity in physics is constant during free fall?})

\writeline
\writeline
\writeline
\end{tcolorbox}

\end{document}
